\section{Web API}
\label{chap:webApi}

Le specifiche complete delle API sono state documentate secondo lo standard \textbf{OpenAPI 3.1.0} e sono consultabili su \texttt{\href{https://wattguard.me/api/v1/docs}{https://wattguard.me/api/v1/docs}}
oppure su \href{https://raw.githubusercontent.com/alterax05/wattguard-ing-sw/refs/heads/main/docs/api-doc.yaml}{GitHub}

\paragraph{Progettazione e note di design: }
L'API segue i principi \emph{RESTful}, organizzando le risorse in modo gerarchico sotto il prefisso \texttt{/api/v1}. 
Si è deciso inoltre di definire gli schemi di input/output con la libreria \texttt{zod}, integrandoli nel documento OpenAPI tramite la libreria hono-openapi.
In questo modo, si riduce il rischio che la documentazione non rispecchi il comportamento effettivo del codice.
Si è inoltre deciso di utilizzare il paradigma "\texttt{API Link}" per permettere una navigazione più semplice e intuitiva delle risorse fornite dall'API.

\inputminted[breaklines, breakanywhere]{yaml}{chapter/api-doc.yaml}